\subsection{Explorador cartográfico y catálogo de alojamientos}
\label{subsec:explorador_alojamientos}

El explorador cartográfico es la herramienta principal mediante la cual los estudiantes descubren y comparan opciones de alojamiento residencial. Diseñada bajo un patrón de panel dividido (*Split-Screen Layout*), la interfaz combina la potencia visual de un mapa interactivo con la riqueza descriptiva de un catálogo de tarjetas inmobiliarias.

\subsubsection{Construcción del visor interactivo en frontend}

En el cliente web, esta pantalla se articula en el componente MapSearchPage.tsx, apoyado en la suite geoespacial de \textbf{Leaflet}, \textbf{React-Leaflet} y \textbf{Supercluster}:

\begin{itemize}
    \item \textbf{Disposición dividida y desplazamiento independiente}: La pantalla se divide en dos columnas ergonómicas. En el lado izquierdo (o panel superior en dispositivos móviles), se presenta el mapa cartográfico interactivo que permanece anclado a la altura de la ventana; en el lado derecho, se despliega la rejilla con las tarjetas de los inmuebles disponibles, dotada de desplazamiento vertical aislado (overflow-y: auto). Al pasar el cursor sobre una tarjeta, el marcador correspondiente en el mapa se ilumina automáticamente mediante una transición visual.
    \item \textbf{Agrupación por conglomerados espaciales (\emph{Supercluster})}: Para garantizar un rendimiento óptimo de 60 cuadros por segundo al manipular el mapa, los puntos de alojamiento no se renderizan individualmente si su densidad es alta. La biblioteca Supercluster calcula en memoria los grupos de marcadores según el nivel de zoom actual y los límites geográficos visibles (*Bounding Box*), mostrando insignias numéricas que se disuelven dinámicamente al acercar el mapa.
    \item \textbf{Filtrado reactivo con retardo controlado (\emph{Debounce})}: La barra de herramientas permite filtrar viviendas por ciudad, rango de precios, tipo de alquiler (habitación privada o piso completo), número de baños y comodidades indispensables (\emph{amenities} como calefacción, ascensor, aire acondicionado o conexión a internet). Para evitar saturar el servidor con peticiones innecesarias mientras el usuario escribe en la barra de búsqueda o ajusta el control deslizante de precio, se aplica una técnica de retardo (*debounce*) de 300~ms mediante lodash.debounce.
\end{itemize}

\subsubsection{Motor de búsqueda geoespacial y paginación en backend}

En el backend, las peticiones de búsqueda convergen en el método search() de AccommodationController:

\begin{enumerate}
    \item \textbf{Filtrado dinámico con especificaciones JPA}: Se implementó un constructor de consultas basado en \textbf{Spring Data JPA Specifications} (JpaSpecificationExecutor). Este patrón compone dinámicamente predicados SQL en función de los parámetros recibidos en la consulta HTTP (city, minPrice, maxPrice, amenities), evitando la concatenación manual de cadenas y blindando el endpoint contra inyecciones SQL.
    \item \textbf{Filtrado por cuadrante geográfico}: Cuando el usuario navega o desplaza el mapa, el cliente envía las coordenadas de las esquinas del área visible (northEastLat, northEastLng, southWestLat, southWestLng). La consulta en base de datos aplica restricciones espaciales de acotamiento (BETWEEN), retornando únicamente las viviendas que se encuentran dentro del cuadrante proyectado.
    \item \textbf{Paginación y proyección optimizada}: Para minimizar el ancho de banda transferido, la respuesta no incluye la entidad completa con todas sus fotos y relaciones, sino un DTO ligero (ListingCardDTO) que contiene únicamente el título, precio mensual, coordenadas y la fotografía de portada principal.
\end{enumerate}

